At the near site with 30\% capacity, the median number of distinct rescued people across runs is 294.0, and median peak shelter occupancy is 292.0 of 292 places. We calculate each run's mean stranding duration among people admitted to shelter; the median of those means is 8.72 h. The shelter is full at a median of 80.0 hourly observation points per run. At the far site with 30\% capacity, the median number of distinct rescued people across runs is 306.0, and median peak shelter occupancy is 291.0 of 292 places. We calculate each run's mean stranding duration among people admitted to shelter; the median of those means is 36.89 h. The shelter is full at a median of 0.0 hourly observation points per run. Stranding durations describe admitted people only; they exclude people who were never rescued and do not separately measure queue waits or travel times. Figure~\ref{fig:process} shows the hourly rescue, occupancy and mortality series for both sites. Slow transport can delay admissions early in the event, while lack of shelter space can prevent admissions later. The total number admitted cannot show which constraint limits rescue at each hour.
